\documentclass[11pt]{article}

\usepackage[a4paper,margin=1in]{geometry}
\usepackage{amsmath,amssymb,amsthm}
\usepackage{booktabs}
\usepackage{hyperref}
\usepackage{xcolor}
\usepackage{microtype}

\newcommand{\Vsix}{V_{600}}
\newcommand{\Vtwentyfour}{V_{24}}
\newcommand{\Lsix}{L_{\Vsix}}
\newcommand{\Loc}{L_{C}}
\newcommand{\Esix}{E_{\Vsix}}
\newcommand{\Elift}{E_{\mathrm{lifted}}}
\newcommand{\Eres}{E_{\mathrm{residual}}}
\newcommand{\Phigold}{\varphi}
\newcommand{\twoI}{2I}
\newcommand{\twoT}{2T}

\theoremstyle{plain}
\newtheorem{theorem}{Theorem}
\newtheorem{proposition}[theorem]{Proposition}
\newtheorem{lemma}[theorem]{Lemma}
\newtheorem{corollary}[theorem]{Corollary}
\theoremstyle{definition}
\newtheorem{remark}[theorem]{Remark}

% Avoid TeX-inserted (discretionary) hyphenated line breaks: these get
% rendered as bad characters in some PDF viewers / extractors.  Keep
% breakable points at existing hyphens (e.g. "v600-programme") since those
% are visible glyphs anyway.  Allow TeX to stretch the line if needed.
\hyphenpenalty=10000
\exhyphenpenalty=100
\setlength{\emergencystretch}{3em}

\title{The 24--600 Spectral Bridge\\
\large A selective $\lambda{=}12$ eigenspace embedding from five 24-cell cosets into the 600-cell}
\author{%
  Lee Smart\\[4pt]
  \textit{Institute of Vibrational Field Dynamics}\\[2pt]
  \texttt{contact@vibrationalfielddynamics.org}\\[2pt]
  \texttt{@vfd\_org}%
}
\date{May 2026}

\begin{document}
\maketitle

\begin{center}
\fbox{\parbox{0.92\textwidth}{\small
\textbf{Status:} Reproducible mathematical note / exact computational
certificate.  Not peer-reviewed.

\medskip
\textbf{Reproducibility.} The formal claims of this note ---
Theorem~\ref{thm:lift} (the $\lambda{=}12$ lift identity); the
$\dim \Esix(\lambda)$ entries for $\lambda\in\{0,4,8,10,12\}$ used
in the selectivity table of Section~\ref{sec:selectivity};
Lemma~\ref{lem:minusone} ($\{\pm 1\}$-triviality);
Lemma~\ref{lem:invariance} ($\twoI$-invariance of $\Elift, \Eres$);
Theorem~\ref{thm:chars} (integer $A_5$ class characters); and
Corollary~\ref{cor:A5} (irrep multiplicities) --- are all certified
by exact rational arithmetic in the accompanying repository (integer
Laplacians, $\mathbb{Q}$-rational nullspace bases, exact zero
residuals; the $A_5$ character projection uses exact
$\mathbb{Q}(\sqrt 5)$ character-table entries).  The numerical
quantities reported in Remark~\ref{rmk:floatcheck} and the
``max-residual'' column of the selectivity table are floating-point
cross-checks only.  One-command reproduction and the frozen output
artifacts are described in Section~\ref{sec:repro}.
}}
\end{center}

\bigskip

\begin{abstract}
The 600-cell is the regular 4-polytope whose 120 vertices form the
binary icosahedral group $\twoI$.  The binary tetrahedral group
$\twoT\subset\twoI$ has index 5; its right cosets partition $\Vsix$
into five disjoint copies of the 24-cell.  Each coset carries a
$d{=}1$ (nearest-neighbour) graph isomorphic to the 24-cell edge
graph (8-regular); on the full vertex set the $d{=}1$ graph is the
600-cell edge graph (12-regular).

We prove that the 2-dimensional $\lambda{=}12$ eigenspace of each
local 24-cell Laplacian zero-extends into the 25-dimensional
$\lambda{=}12$ eigenspace of the full 600-cell Laplacian.  The
identity is established \emph{exactly over} $\mathbb{Q}$: from the
local $\mathbb{Q}$-rational nullspace of each $L_C - 12 I$ (an
integer matrix), every zero-extended basis vector is verified to
lie in the kernel of $\Lsix - 12 I$ by exact rational matrix
multiplication.  A floating-point cross-check independently gives
residual $\le 6.2\times 10^{-15}$.

Among the five local 24-cell Laplacian eigenvalues
$\{0,4,8,10,12\}$, $\lambda{=}12$ is the only one that gives a full
zero-extension into a corresponding global 600-cell eigenspace.  The
$\lambda{=}0$ case contributes exactly one direction (the global
constant, reconstructed from the symmetric sum of the five coset
indicators) of the 5-dimensional lift image; the local eigenvalues
$\lambda\in\{4,8,10\}$ are not eigenvalues of the 600-cell Laplacian
at all, so the question is vacuous there.

Under the induced $A_5 = \twoI/\{\pm1\}$ action on the five cosets,
$\Esix(12) = \Elift \oplus \Eres$ with $\Elift\cong 2\!\cdot\!Y_5$,
$\Eres\cong 3\!\cdot\!Y_5$, $\Esix(12)\cong 5\!\cdot\!Y_5$, where
$Y_5$ is the 5-dimensional irreducible representation of $A_5$.
\end{abstract}

\section{Introduction}

The 600-cell is the regular 4-polytope $\{3,3,5\}$ with 120 vertices,
720 edges and icosahedral symmetry.  Its vertex set is naturally
identified with the binary icosahedral group $\twoI\subset
\mathrm{Sp}(1)$.  The 24-cell is the regular 4-polytope $\{3,4,3\}$
with 24 vertices, 96 edges and $F_4$-symmetry; its vertex set is the
binary tetrahedral group $\twoT$.

It is classical that $\twoT$ has index 5 in $\twoI$ and that the
five cosets of $\twoT$ in $\twoI$ form five inscribed 24-cells whose
vertex sets partition $\Vsix$~\cite{Coxeter1973, ConwaySloane1999}.
An explicit machine-verified treatment within the present programme
is~\cite{SmartXXXV}: the Schl\"afli decomposition is proved in
section~\texttt{sec:bio:schlafli}
(Theorem~\texttt{thm:schlafli}; supporting script
\texttt{papers/}\allowbreak\texttt{cascade-derivation/}\allowbreak\texttt{scripts/}\allowbreak\texttt{verify\_schlafli\_via\_2T.py}),
and the $A_5$ origin of the coset action is proved in
section~\texttt{sec:gr:lorentz}
(Theorem~\texttt{thm:A5action}; script
\texttt{papers/}\allowbreak\texttt{cascade-derivation/}\allowbreak\texttt{scripts/}\allowbreak\texttt{coset\_action\_2I\_on\_5\_skeletons.py}).

The question this note answers is sharper: \emph{does the spectral
structure of each local 24-cell graph embed exactly into the spectral
structure of the global 600-cell graph?}  The answer is ``yes, but
only at $\lambda{=}12$ in full''.  Full zero-extension to a global
$\lambda$-eigenspace occurs only at $\lambda{=}12$; at $\lambda{=}0$
exactly one symmetric combination (the sum of the five coset
indicators, equal to the global constant) lands in the global
eigenspace, while at every other tested $\lambda$ the lift image
intersects the global eigenspace only at $0$.

\paragraph{Why the comparison is non-trivial.}
The local 24-cell Laplacian $\Loc$ and the global 600-cell Laplacian
$\Lsix$ are built from \emph{different} graphs: $\Loc$ uses the
24-cell edge graph on the 24 vertices of a single coset (edges at
squared distance $1$), while $\Lsix$ uses the 600-cell edge graph on
all 120 vertices (edges at squared distance $(3-\sqrt5)/2 < 1$).
These graphs are not subgraphs of one another; in fact \emph{no} edge
of $\Loc$ is an edge of $\Lsix$ (the local 24-cell edge length is the
next distance shell up of $\Vsix$).  Consequently, for a function $v$
supported on a coset $C_i$, the value $(\Lsix v)(u)$ at a vertex
$u\notin C_i$ is in general non-zero --- it is minus the sum of $v$
over the global neighbours of $u$ that happen to lie in $C_i$ --- so
the zero-extension of $v$ is generically \emph{not} a
$\Lsix$-eigenvector.  The content of Theorem~\ref{thm:lift} is that
this off-coset coupling cancels \emph{identically} on the local
$\lambda{=}12$ eigenspace --- a finite-rank linear condition that
turns out to be satisfied --- and the content of the selectivity
table is that no other local eigenspace satisfies it.  This is what
makes the $\lambda{=}12$ channel a genuine local-to-global bridge
rather than a trivial restriction-and-extension.

\section{Construction}

\subsection{$\Vsix$ as $\twoI$}

We work in exact $\mathbb{Q}(\sqrt5)$ arithmetic.  The 120 vertices
of $\Vsix$ are:
\begin{itemize}
  \item 8 vertices $\pm e_i$, $i=0,\dots,3$;
  \item 16 vertices $(\pm 1,\pm 1,\pm 1,\pm 1)/2$;
  \item 96 vertices: even permutations of
        $(0,\pm 1/(2\Phigold),\pm 1/2,\pm \Phigold/2)$,
        $\Phigold=(1+\sqrt5)/2$.
\end{itemize}
These 120 unit quaternions form $\twoI$ under quaternion
multiplication; the construction (exact $\mathbb{Q}(\sqrt5)$
arithmetic, all 120 vertices) lives in the package
\texttt{vfd\_v600}~\cite{Vsixprog}.  Closure of the vertex set under
multiplication is exercised end-to-end by the $A_5$ certificate of
Section~\ref{sec:a5} (which enumerates every $h\in\twoI$ and looks up
$v h\in\twoI$ for every $v\in\twoI$); any non-closure would manifest
as a missing index lookup and abort the certificate.

\subsection{$\Vtwentyfour$ as $\twoT$ = $\sigma$-fixed subset}

The Galois twist $\sigma\colon\sqrt5\mapsto-\sqrt5$ acts component-wise
on $\mathbb{Q}(\sqrt5)$.  Its fixed subset in $\Vsix$ is the
24-element binary tetrahedral group $\Vtwentyfour = \twoT$: the 8
unit vectors $\pm e_i$ together with the 16 half-integer vectors.
Verification is in~\cite{Vsixprog} (\texttt{vfd\_v600.sigma}).

\subsection{Right cosets of $\twoT$ in $\twoI$ and the $A_5$ action}\label{sec:cosets}

\paragraph{Convention.}
We work with the \emph{right cosets} $\twoT\!\cdot\! v$ in $\twoI$;
equivalently, with the orbits of left multiplication by $\twoT$ on
$\twoI$.  Left and right cosets give different partitions of $\twoI$
(since $\twoT$ is not normal in $\twoI$), but both partition the
120-vertex set into five sets of 24, both contain $\twoT$ itself as
the identity coset, and both produce isomorphic 24-cell edge graphs
on each block by isometry.  The right-coset convention is the one
used throughout the keystone-artifact code; the left-coset version
is treated in~\cite{SmartXXXV} (section~\texttt{sec:bio:schlafli}).
The two conventions are related by the involution $g\mapsto g^{-1}$
of $\twoI$, which is a bijection between left and right cosets
fixing the identity coset $\twoT$; hence the spectral-bridge claims
of this note transfer between conventions.

\paragraph{Cosets.}
The five right cosets of $\twoT$ in $\twoI$ are
\begin{itemize}
  \item one coset $C_0 = \twoT$ itself (the $\sigma$-fixed orbit);
  \item four cosets $C_1,\dots,C_4$: the four $\twoT$-orbits on the
        96 $\sigma$-mobile vertices.
\end{itemize}
Each $|C_i| = 24$, and $\bigsqcup C_i = \Vsix$~\cite{SmartXXXV,
Vsixprog}.

\paragraph{Action on cosets.}
$\twoI$ acts on the five right cosets by right multiplication:
$(\twoT\!\cdot\! v)\cdot g = \twoT\!\cdot\!(vg)$.  The kernel of this
\emph{coset-level} action contains $\{\pm1\}$ because $-1 \in \twoT$
(since $\twoT$ contains $\pm e_0 = \pm1$ as one of its 24 elements),
so for any right coset $\twoT v$ we have $\twoT v\cdot(-1) =
\twoT(-v) = (-\twoT)\,v = \twoT v$.  By the parallel left-coset
result in~\cite{SmartXXXV} (proved at all 120 elements of $\twoI$ by
\texttt{papers/cascade-derivation/scripts/coset\_action\_2I\_on\_5\_skeletons.py})
combined with the involution $g \mapsto g^{-1}$ which exchanges
left- and right-coset families, the kernel of this action is exactly
$\{\pm 1\}$ and the image on five letters is exactly
$A_5 = \twoI/\{\pm 1\}$, acting transitively.

\paragraph{Vertex action vs.\ coset action.}
By contrast, right multiplication on \emph{vertices} of $\Vsix$ has
trivial kernel: $v \cdot g = v$ for some $v$ forces $g = 1$.  When
we use the term ``$A_5$ action'' on a subspace $E\subset
\mathbb{R}^{\Vsix}$ below, we mean the action induced on $E$ by
right multiplication of vertices, restricted to subspaces on which
$-1$ acts trivially (so that the linear action descends through
$\{\pm 1\}$ to $A_5$).  Whether $-1$ acts trivially on a given
$E\subset\mathbb{R}^{\Vsix}$ is a non-trivial fact specific to $E$;
we prove it for $E = \Esix(12)$ in
Section~\ref{sec:a5} (Lemma~\ref{lem:minusone}).

\medskip
\noindent
\begin{tabular}{lrrl}
\toprule
Object & $|V|$ & $|E|$ & Graph (= shortest non-zero distance shell)\\
\midrule
$\Vsix$ & 120 & 720 & shortest-$\Vsix$ shell, squared distance $(3-\sqrt5)/2$, 12-regular\\
$\twoT$-coset (any) & 24 & 96 & shortest intra-coset shell, squared distance $1$, 8-regular\\
\bottomrule
\end{tabular}

\smallskip\noindent
The two distance shells are distinct: \emph{no} pair within a single
coset has separation $(3-\sqrt5)/2$, and conversely no $\Vsix$-edge
joins two vertices of the same coset.  Throughout this note, ``$d{=}1$
graph'' is shorthand for ``shell-1 (nearest-neighbour) graph at the
shortest non-zero distance of whichever vertex set we are working on'';
the squared distances are different for $\Vsix$ and a coset, as the
table makes explicit.

\subsection{The two $d{=}1$ shells}\label{sec:two-shells}

Inside $\mathbb{Q}(\sqrt5)$, there are 8 distinct non-zero squared
distances between pairs in $\Vsix$.  The shortest is
\[
  d^2_{V_{600}} \;=\; \tfrac{3-\sqrt5}{2} \;\approx\; 0.382,
\]
which defines the 600-cell edge graph (12-regular, 720 edges).
Inside each right coset $C_i$, the shortest non-zero squared distance
between members is the next shell up:
\[
  d^2_C \;=\; 1,
\]
which on $C_i$ realises the standard 24-cell edge graph (8-regular,
96 edges).  The two distances are distinct: \emph{no} pair within a
single coset has separation $d^2_{V_{600}}$.

These two distinct shells underlie the two Laplacians compared in
this note.

\subsection{Laplacians}\label{sec:laplacians}

We use the unweighted combinatorial Laplacian $L = D - A$.  Both
graphs are regular and have integer 0/1 adjacency matrices and
integer regular degree, so the Laplacians are integer matrices.
Numerical diagonalisation (the local spectrum agrees with the
floating-point computation in \cite{SmartV},
\texttt{papers/paper-v/scripts/run\_polytope\_uniqueness\_narrow.py};
every integer-eigenvalue multiplicity in both spectra below is
additionally certified by exact $\mathbb{Q}$-rational SymPy
nullspace of $\Lsix - \lambda I$ and $\Loc - \lambda I$ at the
specific integer $\lambda$ values shown in the spectra, via
\texttt{closure\_transform\_engine.keystone.exact\_integer\_spectrum\_check}):
\[
  \mathrm{spec}\,\Loc \;=\; \{0^{\,1},\,4^{\,4},\,8^{\,9},\,10^{\,8},\,12^{\,2}\},
\]
\[
  \mathrm{spec}\,\Lsix \;=\;
   \{0^{\,1},\,(\approx\!2.29)^{\,4},\,(\approx\!5.53)^{\,9},\,9^{\,16},\,
     12^{\,25},\,14^{\,36},\,(\approx\!14.47)^{\,9},\,15^{\,16},\,(\approx\!15.71)^{\,4}\},
\]
in multiset notation $\lambda^{\,m}$; the four non-integer
eigenvalues take values in $\mathbb{Q}(\sqrt 5)$.  The integer
eigenvalues of $\Lsix$ (including the multiplicity
$\dim\Esix(12) = 25$ at $\lambda = 12$ used in
Corollary~\ref{cor:dim} and Section~\ref{sec:a5}) coincide with
the $\mathbb{Q}$-rational kernel dimensions computed by the SymPy
nullspace routine in
\texttt{closure\_transform\_engine.keystone}, so $\dim\Esix(12)$
is itself an exact rational-arithmetic spectral fact.

\section{The $\lambda{=}12$ lift identity}\label{sec:lift}

For a coset $C_i$, define the zero-extension map
\[
  \mathrm{lift}_i\colon \mathbb{R}^{24}\to\mathbb{R}^{120},\qquad
  (\mathrm{lift}_i u)(v) = \begin{cases} u(v) & v\in C_i,\\
                                          0    & v\notin C_i.\end{cases}
\]

\begin{theorem}[$\lambda{=}12$ lift identity]\label{thm:lift}
For every coset $C_i$,
\[
  \mathrm{lift}_i\bigl(E_{C_i}(12)\bigr) \;\subset\; \Esix(12).
\]
\end{theorem}

\begin{proof}[Proof by exact rational certificate]
Since the local Laplacian $\Loc$ has integer entries and 12 is an
integer eigenvalue, $E_{C_i}(12) = \ker(\Loc - 12I)$ is a
$\mathbb{Q}$-rational subspace of $\mathbb{Q}^{24}$, and $\Lsix$
likewise has integer entries.  The accompanying repository produces
a checkable certificate consisting of:
(i) the integer matrices $\Loc - 12I$ and $\Lsix - 12I$;
(ii) a $\mathbb{Q}$-rational basis $\{v_1, v_2\}$ of each
$E_{C_i}(12) = \ker(\Loc - 12I)$, obtained by exact Gaussian
elimination over $\mathbb{Q}$;
(iii) for each $v_k$, the vector $(\Lsix - 12I)\,\mathrm{lift}_i(v_k)
\in \mathbb{Q}^{120}$, computed by exact rational matrix
multiplication, with every one of its 120 coordinates equal to $0$.
Item (iii) is the certificate of $\Lsix\,\mathrm{lift}_i(v_k) =
12\,\mathrm{lift}_i(v_k)$ in $\mathbb{Q}^{120}$, for both $k=1,2$
and all five cosets $i$; no floating-point arithmetic enters.  The
certificate is regenerated and re-checked by
\texttt{closure\_transform\_engine.keystone.exact\_certify\_lambda12\_lift}
and asserted by the test suite.
\end{proof}

\begin{remark}[Floating-point cross-check]\label{rmk:floatcheck}
Independent floating-point diagonalisation of $\Lsix$ and each $\Loc$
via \texttt{numpy.linalg.eigh}, followed by zero-extension of each
basis vector $v$ of $E_{C_i}(12)$ and evaluation of
$\|\Lsix v_{\mathrm{lifted}} - 12 v_{\mathrm{lifted}}\|$, gives
maximum residual $\le 6.2\times10^{-15}$ across all 10 lifted basis
vectors.  This is consistent with Theorem~\ref{thm:lift} at the
level of numerical precision.
\end{remark}

\begin{corollary}\label{cor:dim}
Setting $\Elift := \bigoplus_i \mathrm{lift}_i(E_{C_i}(12))$, the
five lifted local eigenspaces are linearly independent (their
supports are the disjoint cosets), so
$\dim \Elift = 5\cdot 2 = 10$.  Since $\dim \Esix(12) = 25$, the
subspace $\Elift\subset\Esix(12)$ is proper.
\end{corollary}

\section{Selectivity}\label{sec:selectivity}

We test whether the analogue of Theorem~\ref{thm:lift} holds for
other eigenvalues $\lambda$ of $\Loc$.  The natural measure is the
\emph{intersection dimension}
\[
  \mathfrak{s}(\lambda) \;:=\; \dim\Bigl(\;
     \mathrm{Im}\bigl(\mathrm{lift}\bigr)\big|_{\lambda}
     \;\cap\;\Esix(\lambda)\;\Bigr),
\]
where $\mathrm{Im}(\mathrm{lift})\big|_{\lambda} = \sum_i
\mathrm{lift}_i\bigl(E_{C_i}(\lambda)\bigr)\subset\mathbb{R}^{120}$.
A ``full'' lift would have $\mathfrak{s}(\lambda)$ equal to the
total lift-image dimension; a generic eigenvalue has
$\mathfrak{s}(\lambda)=0$.

\begin{center}
\begin{tabular}{rrrrl}
\toprule
local $\lambda$ & lift-image dim & $\dim \Esix(\lambda)$ &
   $\mathfrak{s}(\lambda)$ & verdict \\
\midrule
0  & 5  & 1  & 1 & PARTIAL — only global constant \\
4  & 20 & 0  & 0 & $\Esix(\lambda) = \{0\}$ \\
8  & 45 & 0  & 0 & $\Esix(\lambda) = \{0\}$ \\
10 & 40 & 0  & 0 & $\Esix(\lambda) = \{0\}$ \\
12 & 10 & 25 & 10 & FULL \\
\bottomrule
\end{tabular}
\end{center}

The PARTIAL row at $\lambda{=}0$ is the substantive control: here
$\Esix(0)$ is one-dimensional (the global constant), and one might
ask whether the lift hits it.  It does --- but only through the
symmetric combination: $\mathbf{1}_{V_{600}}$ is the sum of the five
coset indicators, each a local $\lambda{=}0$ eigenvector, so exactly
one direction of the 5-dimensional lift image lands in $\Esix(0)$.
No single-coset $\lambda{=}0$ lift lies in $\Esix(0)$; only the
$A_5$-invariant combination does.  This is the sense in which even
the one eigenvalue that the lift ``partly'' reaches, it reaches only
in the symmetry-forced way.

For $\lambda\in\{4,8,10\}$ the control is weaker: these are simply
not eigenvalues of $\Lsix$ at all.  The exact $\mathbb{Q}$-rational
sympy nullspace of $\Lsix - \lambda I$ has dimension $0$ for each of
$\lambda\in\{4,8,10\}$ (certified by
\texttt{exact\_integer\_spectrum\_check}), so $\Esix(\lambda) = \{0\}$
and $\mathfrak{s}(\lambda) = 0$ holds for trivial dimensional reasons.
We list these rows for completeness --- and the floating-point lift
residuals at these $\lambda$ ($3.16$ to $9.38$, far above machine
precision) confirm the lifted vectors are nowhere near an
approximate $\lambda$-eigenfunction of $\Lsix$ --- but we do not
claim they constitute a strong falsification test in the way the
$\lambda{=}0$ row does.

Hence among the five eigenvalues of $\Loc$, only $\lambda{=}12$
satisfies the lift identity in full; at $\lambda{=}0$ the lift is
partial in the single symmetry-forced way; and at the remaining
$\lambda$ the question is empty for spectral reasons.

\section{Direct-sum decomposition and $A_5$ structure}\label{sec:a5}

Inside $\Esix(12)$, the lifted subspace $\Elift$ (dim 10) has an
orthogonal complement
\[
  \Eres \;:=\; \Esix(12)\;\ominus\;\Elift,\qquad \dim\Eres = 15.
\]
We call $\Eres$ the \emph{residual subspace}; it consists of those
global $\lambda{=}12$ eigenvectors that are orthogonal to every
zero-extension of a local $\lambda{=}12$ eigenfunction.

\subsection{Descent of right multiplication to $A_5$ on $\Esix(12)$}

For every $h\in\twoI$ let $P_h\in\mathrm{Perm}(\Vsix)$ be the
120-permutation matrix induced by right multiplication
$v\mapsto v h$ on the vertex set.  Right multiplication is an
isometry of $\Vsix$, so $P_h$ commutes with $\Lsix$ and restricts to
a linear action on $\Esix(12)$.  Right multiplication on
\emph{vertices} has trivial kernel; descent to an $A_5$ action on a
subspace $E\subset\mathbb{R}^{\Vsix}$ is the separate question of
whether $-1$ acts trivially on $E$.

\begin{lemma}[$\{\pm 1\}$ acts trivially on $\Esix(12)$]\label{lem:minusone}
$P_{-1}$ restricted to $\Esix(12)$ is the identity.
\end{lemma}

\begin{proof}[Proof by exact rational computation]
Compute the $\mathbb{Q}$-rational nullspace
$B_{\mathrm{global}}\in\mathbb{Q}^{120\times 25}$ of $\Lsix - 12 I$
(an integer matrix), giving a rational basis for $\Esix(12)$.  For
each column $v$ of $B_{\mathrm{global}}$, apply the integer
permutation matrix $P_{-1}$ exactly and verify
$P_{-1}\,v - v = 0 \in \mathbb{Q}^{120}$ componentwise (no
floating point).  Every one of the $25\cdot 120 = 3000$ rational
component checks returns $0$.  Hence $P_{-1} = \mathrm{Id}$ on
the basis, and therefore on $\Esix(12)$.  The computation is
reproduced by
\texttt{closure\_transform\_engine.keystone.exact\_a5\_characters}
(field \texttt{minus\_one\_acts\_trivially}).
\end{proof}

By Lemma~\ref{lem:minusone}, the action of $\twoI$ on $\Esix(12)$
factors through $A_5 = \twoI/\{\pm 1\}$.  To restrict this action to
the summands $\Elift$ and $\Eres$ we need $\twoI$-invariance of each.

\begin{lemma}[$\Elift$ and $\Eres$ are $\twoI$-invariant]\label{lem:invariance}
For every $h\in\twoI$, $P_h \Elift = \Elift$ and $P_h \Eres = \Eres$.
\end{lemma}

\begin{proof}
Right multiplication by $h\in\twoI$ permutes the right cosets:
$P_h$ sends $C_i$ bijectively to $C_{\sigma_h(i)}$, where
$\sigma_h\in S_5$ is the induced coset permutation
(Section~\ref{sec:cosets}).  Restricted to a coset, $P_h\!\restriction_{C_i}\colon
C_i \to C_{\sigma_h(i)}$ is an isometry between two metric spaces
that are each isomorphic to the 24-cell.  It therefore intertwines
the two local shell-1 Laplacians: $L_{C_{\sigma_h(i)}} (P_h v) =
P_h (L_{C_i} v)$ for $v\in\mathbb{R}^{C_i}$.  Consequently $P_h$
maps $E_{C_i}(12)$ to $E_{C_{\sigma_h(i)}}(12)$, and by linearity of
zero-extension, $P_h$ maps $\mathrm{lift}_i(E_{C_i}(12))$ to
$\mathrm{lift}_{\sigma_h(i)}(E_{C_{\sigma_h(i)}}(12))$.  Summing
over $i$ gives $P_h \Elift \subseteq \Elift$; equality follows from
invertibility of $P_h$.  $\Eres$ is the orthogonal complement of
$\Elift$ in $\Esix(12)$, and $P_h$ is orthogonal (a permutation
matrix), so $P_h$ also preserves $\Eres$.

\smallskip
\noindent
Independent of this prose proof, the repository carries an exact
$\mathbb{Q}$-rational certificate: for every $h\in\twoI$ and for
$B\in\{B_{\Elift},\,B_{\Eres}\}$, sympy computes
$P_h B - B\,(B^\top B)^{-1} B^\top(P_h B)$ over $\mathbb{Q}$ and
verifies that every component is exactly $0$.  This certifies
$P_h B \subseteq \mathrm{col}(B)$, hence $P_h$ preserves $\Elift$
and $\Eres$, with no reference to the prose argument.  Reproduced by
the \texttt{lifted\_2I\_invariant} and \texttt{residual\_2I\_invariant}
fields of \texttt{exact\_a5\_characters}.
\end{proof}

By Lemmas \ref{lem:minusone}--\ref{lem:invariance}, the $A_5$ action
on $\Esix(12)$ restricts to each of $\Elift$ and $\Eres$.

\subsection{Exact $A_5$ class characters}

Using rational bases of $\Elift$, $\Eres$, and $\Esix(12)$
constructed by sympy nullspace, we compute the trace of $P_h$ on
each subspace by exact rational matrix multiplication for every
$h\in\twoI$, and average by $A_5$ conjugacy class.  The five
$A_5$ conjugacy classes 1A, 2A, 3A, 5A, 5B are obtained by
explicit orbit enumeration: we first enumerate the 60 distinct
5-letter permutations $\pi(h)$ for $h\in\twoI$ (the image of the
coset action), then partition them into orbits under conjugation
$g\mapsto \pi\,g\,\pi^{-1}$ by all $\pi\in\pi(\twoI)$.  The
non-trivial cycle structures $(1,1,1,1,1)$, $(1,2,2)$ and $(1,1,3)$
each give a single $A_5$ class (sizes 1, 15, 20); the cycle
structure $(5)$ splits into two $A_5$ orbits of size 12 each, which
we label 5A and 5B.  Class sizes $(1, 15, 20, 12, 12)$ match the
standard $A_5$ character table.

\begin{theorem}[Exact $A_5$ class characters]\label{thm:chars}
The $A_5$ class characters of $\Esix(12)$, $\Elift$, and $\Eres$
are constant on each conjugacy class and take exact integer values:

\smallskip
\begin{tabular}{lrrrrr}
\toprule
subspace & 1A & 2A & 3A & 5A & 5B \\
\midrule
$\Esix(12)$ & 25 & 5  & $-5$ & 0 & 0 \\
$\Elift$    & 10 & 2  & $-2$ & 0 & 0 \\
$\Eres$     & 15 & 3  & $-3$ & 0 & 0 \\
\bottomrule
\end{tabular}
\smallskip

Each entry is computed in $\mathbb{Q}$ (no floating point).
\end{theorem}

Projecting the exact characters of Theorem~\ref{thm:chars} against
the $A_5$ character table (irreps of dimensions $1, 3, 3, 4, 5$;
see \cite{ConwayAtlas1985}) gives:

\begin{center}
\begin{tabular}{lrrrr}
\toprule
irrep & dim & in $\Elift$ & in $\Eres$ & in $\Esix(12)$\\
\midrule
$Y_1$  & 1 & 0 & 0 & 0 \\
$Y_3$  & 3 & 0 & 0 & 0 \\
$Y_3'$ & 3 & 0 & 0 & 0 \\
$Y_4$  & 4 & 0 & 0 & 0 \\
$Y_5$  & 5 & 2 & 3 & 5 \\
\bottomrule
\end{tabular}
\end{center}

\begin{corollary}[$A_5$ irrep decomposition]\label{cor:A5}
With $Y_5$ the 5-dimensional irreducible representation of $A_5$,
\[
  \Elift \cong 2\!\cdot\! Y_5,\qquad
  \Eres  \cong 3\!\cdot\! Y_5,\qquad
  \Esix(12) \cong 5\!\cdot\! Y_5,
\]
with every other complex $A_5$ irrep having multiplicity exactly $0$
(the character inner products that compute these multiplicities are
themselves computed in $\mathbb{Q}(\sqrt 5)$ using the exact
character-table entries).
\end{corollary}

The ``$15 = 3\times 5$'' is representation-theoretic — three copies
of the 5-dim $A_5$ irrep — \emph{not} five cosets each contributing
three local modes.

\section{Interpretation and limits}

This is a reproducible spectral-graph result in the 24-cell/600-cell
coset geometry.  It does not by itself derive particle physics,
quantum mechanics, or cosmology.

\smallskip
\noindent
Within the broader closure-geometry programme that motivated this
work, the result may be interpreted as one example of a
``closure-projection channel'': a local shell structure lifting
exactly into a global invariant sector.  That interpretation is
provided here only for context and is not load-bearing for any claim
in this note.

\section{Reproducibility}\label{sec:repro}

\paragraph{One-command reproduction.}
From the repository root:
\begin{verbatim}
python closure_transform_engine/examples/run_wo008_keystone.py
pytest  closure_transform_engine/tests/test_wo008_keystone_artifact.py
\end{verbatim}
The first command (a few minutes on commodity hardware) builds the
600-cell from exact $\mathbb{Q}(\sqrt5)$ icosian quaternions, runs
all certificates, and writes the frozen output artifacts listed
below; the second re-asserts the certificate-level claims listed
below as a test.

\paragraph{Environment.}
The repository pins the toolchain in \texttt{requirements.txt} and
\texttt{pyproject.toml}: Python~$\ge 3.8$ (developed on 3.8.10),
NumPy~$\ge 1.20$ (1.24.4), SymPy~$\ge 1.12$ (1.13.3), SciPy~$\ge 1.7$
(1.10.1).  The formal certificates use only exact $\mathbb{Q}$
arithmetic (SymPy \texttt{Rational}/\texttt{Matrix}); NumPy enters
only the floating-point cross-check of Remark~\ref{rmk:floatcheck}.
The \texttt{wo008\_reproducibility\_log.txt} output records the run
UTC timestamp, elapsed time, Python version, repository path, and
the repository's commit hash (with an
``uncommitted changes present'' marker if the working tree was dirty
at run time), so the artifact set is anchored to a specific source
revision.

\paragraph{Repository contents (proof package).}
The repository is part of the proof package, not an afterthought.  It
contains: the exact construction of $\Vsix$, $\twoT$, the five
cosets, and the $d{=}1$ graph Laplacians; the exact nullspace and
certificate routines for Theorem~\ref{thm:lift},
Lemma~\ref{lem:minusone}, Theorem~\ref{thm:chars}, and the integer
spectrum; the $A_5$ conjugacy-class enumeration and character
projection; the test suite; and the frozen output artifacts
(\texttt{wo008\_summary.json}, \texttt{wo008\_verdict\_table.csv},
\texttt{wo008\_reproducibility\_log.txt}).  The companion scripts cited from~\cite{SmartV} and~\cite{SmartXXXV}
live in the main VFD research repository (paths given in the
bibliography); the self-contained reproduction bundle for this note
vendors only the \texttt{vfd\_v600} icosian package it depends on.
\emph{The claims of this note do not require rerunning Paper~V or
Paper~XXXV.}  Those upstream references justify the broader
programme context and provide an independent treatment of the coset
construction; the standalone bundle reconstructs from scratch the
objects needed for the certificates used here, so a reader can
verify every theorem of this note without access to the main
research repository.

\medskip
\noindent What the runner verifies:
\begin{itemize}
  \item Theorem~\ref{thm:lift} via the $\mathbb{Q}$-rational
        certificate (integer matrices, $\mathbb{Q}$-rational nullspace
        bases, exact zero off-coset components after lift in every one
        of $5\cdot 2 = 10$ local basis vectors).
  \item The floating-point cross-check of
        Remark~\ref{rmk:floatcheck} (max residual
        $\le 10^{-10}$; observed $\approx 6.2\times 10^{-15}$).
  \item Every integer-eigenvalue multiplicity of $\Loc$ and $\Lsix$
        listed in the spectra (exact $\mathbb{Q}$-rational sympy
        nullspace at each integer $\lambda$).
  \item The intersection-dimension table of
        Section~\ref{sec:selectivity}.
  \item Lemma~\ref{lem:minusone} ($-1\in\twoI$ acts trivially on
        $\Esix(12)$) by exact $\mathbb{Q}$-rational basis check.
  \item Lemma~\ref{lem:invariance} ($\twoI$-invariance of $\Elift$
        and $\Eres$): exact $\mathbb{Q}$-rational check of
        $P_h B = B\,(B^\top B)^{-1} B^\top(P_h B)$ for every
        $h\in\twoI$ and both $B=B_{\Elift},\,B_{\Eres}$.
  \item Theorem~\ref{thm:chars} (exact integer $A_5$ class
        characters) and Corollary~\ref{cor:A5} ($A_5$ irrep
        multiplicities) by exact rational trace and character
        projection.
\end{itemize}

Outputs:
\begin{itemize}
  \item \texttt{wo008\_summary.json} — full structured result
        including the certificate's per-coset status.
  \item \texttt{wo008\_verdict\_table.csv}
  \item \texttt{wo008\_reproducibility\_log.txt}
\end{itemize}

\begin{thebibliography}{9}
\bibitem[Coxeter 1973]{Coxeter1973}
H.\,S.\,M.~Coxeter, \emph{Regular Polytopes}, 3rd ed., Dover, 1973.
\bibitem[Conway--Sloane 1999]{ConwaySloane1999}
J.\,H.~Conway and N.\,J.\,A.~Sloane, \emph{Sphere Packings, Lattices and
Groups}, 3rd ed., Springer-Verlag, 1999.
\bibitem[Conway et al.\ 1985]{ConwayAtlas1985}
J.\,H.~Conway, R.\,T.~Curtis, S.\,P.~Norton, R.\,A.~Parker and R.\,A.~Wilson,
\emph{Atlas of Finite Groups}, Oxford University Press, 1985.
\bibitem[Paper~V]{SmartV}
\emph{Polytope-uniqueness theorem and spectra.}  In the main VFD
research repository:
\texttt{papers/paper-v/paper-v.tex}; supporting scripts
\texttt{papers/paper-v/scripts/run\_polytope\_uniqueness\_narrow.py}
and \texttt{papers/paper-v/scripts/run\_rigorous\_verification.py}.
\bibitem[Paper~XXXV]{SmartXXXV}
\emph{Schl\"afli decomposition of the 600-cell and the $A_5$ origin of
the five inscribed 24-cells.} In this repository:
\texttt{papers/paper-xxxv/paper-xxxv.tex}; supporting scripts at
\texttt{papers/cascade-derivation/scripts/verify\_schlafli\_via\_2T.py}
and
\texttt{papers/cascade-derivation/scripts/coset\_action\_2I\_on\_5\_skeletons.py}.
Paper~XXXV proves the left-coset version explicitly; the right-coset
version used here follows from the involution $g\mapsto g^{-1}$ of
$\twoI$, which exchanges the left- and right-coset families and
preserves their common identity coset.
\bibitem[V$_{600}$-programme]{Vsixprog}
\emph{Exact $\mathbb{Q}(\sqrt5)$ icosian arithmetic and
$\twoI$/$\twoT$/$\sigma$-Galois machinery.}  In this repository:
\texttt{papers/v600-programme/lib/vfd\_v600/}.
\end{thebibliography}

\end{document}
